\scanpage{30}
% @ 1–4
\uline{Corollaire} 1 \unsure{Tout|Pour} \deleted{\unsure{tout}} sous-module $M'$ de $M$, il y a un plus petit sous-module \tikz[baseline=(s.base)]{\node[inner sep=1pt](s){\unsure{stble|stbl}};\node[anchor=south east,inner sep=1pt] at ([xshift=4pt,yshift=6pt]s.north west){\inserted{$M'^{\unsure{\vcenter{\hbox{\tikz[x=1ex,y=1ex]{\draw[line width=.11ex,line cap=round,line join=round] (.57,.90) .. controls (.45,.76) and (.24,.53) .. (.15,.35) .. controls (.35,.35) and (.56,.45) .. (.73,.41) .. controls (.64,.2) and (.45,.03) .. (.32,-.14);}}}|\natural|\sharp|\mathrm{s}}}$}};\draw ([xshift=-8pt,yshift=-2pt]s.south west) to[out=80,in=200] ([xshift=1pt,yshift=7pt]s.north west) to[out=20,in=250] ([xshift=8pt,yshift=13pt]s.north west);} de $M$ qui contienne $M'$, et c'est aussi la sous $\check A$-algèbre engendrée. Si $M'$ est de t.f. \deleted{sur} $k$, $M'^{\vcenter{\hbox{\tikz[x=1ex,y=1ex]{\draw[line width=.11ex,line cap=round,line join=round] (.57,.90) .. controls (.45,.76) and (.24,.53) .. (.15,.35) .. controls (.35,.35) and (.56,.45) .. (.73,.41) .. controls (.64,.2) and (.45,.03) .. (.32,-.14);}}}}$ aussi, \deleted{donc} \deleted{\unsure{M}}% end of the crossed group

% @ 5–6
\uline{Corollaire 2} $M$ est limite inductive filtr.\deleted{\unsure{$\to$|t}} de ses sous-modules de t. f. \uline{stbles}, \deleted{\unsure{par}}

% @ 7–8
\deleted{\uline{Corollaire 3}} \deleted{\unsure{M}} \deleted{Appliquons} Considérons $A$ comme bicomodule sur lui-même à g. et à droite, \deleted{plus} donc comme bi-module sur $\check A$,
% @ 9–12
\deleted{Alors pour \unsure{H|tt} $x$, les \unsure{s.}} considérons les sous-bi-modules $B\subset A$. \deleted{On voit que pour} Appliquant Cor 1 et 2 \deleted{les sous-\unsure{bi-mod.}} \deleted{\unsure{$A$|à} $A$ considéré comme $A$} à \tikz[baseline=(s.base)]{\node[inner sep=1pt](s){\(A\otimes A^{\mathrm{o}}\)};\node[anchor=south west,inner sep=1pt](a) at ([xshift=-5pt,yshift=8pt]s.north west){\inserted{la coalgèbre produit ; ($A^{\mathrm{o}}$ l'opposée de $A$)}};\draw ([xshift=-13pt,yshift=-3pt]s.south west) to[out=75,in=190] (a.south west) to[out=5,in=180] ([xshift=3pt,yshift=-2pt]a.south east) to[out=0,in=-90] ([xshift=4pt,yshift=3pt]a.south east);} et au module $A$ pour celle-ci, ou raisonnant directement, on trouve

% @ 13–16
\uline{Cor 3} \deleted{\unsure{Tout} \unsure{bi} \unsure{linéi}} Tout sous-bi-module $B$ de $A$ est contenu dans un plus-petit sous-bi-module $B^{\vcenter{\hbox{\tikz[x=1ex,y=1ex]{\draw[line width=.11ex,line cap=round,line join=round] (.57,.90) .. controls (.45,.76) and (.24,.53) .. (.15,.35) .. controls (.35,.35) and (.56,.45) .. (.73,.41) .. controls (.64,.2) and (.45,.03) .. (.32,-.14);}}}}$ \deleted{\unsure{stable|stble} par} \deleted{\inserted{\unsure{qui est}}} st.bl à g. et à droite (savoir le sous-$\check A$-bi-module engendré), et si $B$ est de t. f. sur $k$, $B^{\vcenter{\hbox{\tikz[x=1ex,y=1ex]{\draw[line width=.11ex,line cap=round,line join=round] (.57,.90) .. controls (.45,.76) and (.24,.53) .. (.15,.35) .. controls (.35,.35) and (.56,.45) .. (.73,.41) .. controls (.64,.2) and (.45,.03) .. (.32,-.14);}}}}$ aussi.
% @ 17–18
Donc $A=\underset{\mathrel{\tikz[baseline=-.3ex,x=1em,y=1em]{\draw[line width=.45pt] (0,0)--(1,.05)--(1,-.10);}}}{\mathrm{lim}}\,A_i$, les $A_i$ étant st.bles à g. et à droite.

% @ 19–23
Or cette st.bilité \tikz[baseline=(s.base)]{\node[inner sep=1pt](s){signifie};\node[anchor=south east,inner sep=1pt](a) at ([xshift=3pt,yshift=6pt]s.north west){\inserted{pour un $B$ donné}};\draw ([xshift=-2pt,yshift=-1pt]s.south west) to[out=90,in=180] ([xshift=2pt,yshift=-2pt]a.south west)--([xshift=-1pt,yshift=-2pt]a.south east) to[out=0,in=-90] ([xshift=2pt,yshift=1pt]a.south east);} que $\unsure{\Pi|\pi}:A\mathrel{\tikz[baseline=-.3ex,x=1em,y=1em]{\draw[line width=.45pt] (0,0)--(1,.05)--(1,-.10);}}A\otimes A$ envoie $B$ dans $(B\otimes A)\unsure{.|,}\cap(A\otimes B)$. Si \deleted{$B$ \unsure{est}} $A/B$ est plat sur $k$ (p ex $k$ un corps) cel signifie aussi que $\Pi$ envoie $B$ dans $B\otimes B$, ou encore qu'il y a une structure de co-algèbre induite. On trouve donc

% @ 24–26
\uline{Cor. 4} Toute coalgèbre \tikz[baseline=(s.base)]{\node[inner sep=1pt](s){sur};\node[anchor=south east,inner sep=1pt] at ([xshift=2pt,yshift=7pt]s.north west){\inserted{$\unsure{A|\wedge}$}};\draw ([xshift=-8pt,yshift=-5pt]s.south west) to[out=90,in=200] ([xshift=-2pt,yshift=5pt]s.north west) to[out=20,in=270] ([xshift=6pt,yshift=11pt]s.north west);} un corps est \tikz[baseline=(s.base)]{\node[inner sep=1pt](s){\(\underset{\mathrel{\tikz[baseline=-.3ex,x=1em,y=1em]{\draw[line width=.45pt] (0,0)--(1,.05)--(1,-.10);}}}{\mathrm{lim}}\)};\node[anchor=south west,inner sep=1pt] at ([xshift=3pt,yshift=6pt]s.north east){\inserted{fil.}};\draw ([xshift=5pt,yshift=-2pt]s.south east) to[out=65,in=225] ([xshift=19pt,yshift=17pt]s.north east);} de coalgèbres de t. f., donc $\check A$ est \tikz[baseline=(s.base)]{\node[inner sep=1pt](s){\(\underset{\leftarrow}{\mathrm{lim}}\)};\node[anchor=south west,inner sep=1pt](a) at ([xshift=-1pt,yshift=7pt]s.north west){\inserted{fil. stricte}};\draw ([xshift=4pt,yshift=-3pt]s.south east) to[out=90,in=190] ([xshift=4pt,yshift=-1pt]a.south east) to[out=10,in=-90] ([xshift=9pt,yshift=6pt]a.south east);} d'algèbres finies sur $k$.
% archive sheet number: 29
